\section{问题定义与方法概览}\label{sec:problem}
\subsection{同值发布为什么需要保留身份}\label{sec:example}
考虑一个双线程、两轮的重复发布实例。线程 A 每轮先写 $x$，发布一次标志 $fa$，再写 $y$，随后再次发布同一标志值；线程 B 对称地先写 $y$，发布 $fb$，再写 $x$，再次发布 $fb$。两者接着读对方的标志并读取本线程先写的数据地址。数据写值均为 $1$，第一、二轮标志分别写 $1,0$；对应等待读分别观测 $0,1$。附录~\ref{app:witness} 列出所有事件和合法完成。与当前查询相关的事件是
\[
\begin{aligned}
 A\text{ 第1轮}:&\quad \cdots\ a_0=\W(fa,1)\ ;\ u=\W(y,1)\ ;\ a_1=\W(fa,1)\ \cdots,\\
 B\text{ 第2轮}:&\quad v=\W(y,1)\ ;\ \cdots\ ;\ r=\R(fa,1)\ ;\ s=\R(y,1).
\end{aligned}
\]
这是一段发布内核的有限观测，用来展示查询语义。第~\ref{sec:protocols}节的循环缓冲与屏障另有可追溯的协议来源。

令 $d=\RF vs$，$e=\RF{a_0}r$，$l=\RF{a_1}r$，$q=\CO uv$。在假设 $d\land l$ 下，有
\[
 u\mathrel{\SB}a_1\mathrel\rf r\mathrel{\SB}s,
\]
所以 $u\mathrel\hb s$；由 WR 一致性和 $s$ 的来源 $v$ 得到 $q$。于是 $d\land l\land\neg q$ 不可满足。若将来源换为早发布 $a_0$，这一因果路径消失，而且 $d\land e\land\neg q$ 确有合法完成；$d\land e\land q$ 也有合法完成。两种标志来源写入相同值，却对另一个地址的 CO 支持产生不同影响。撤销 $l$ 后，$d$ 下两个 CO 方向恢复支持。这里“恢复”有完整见证，而不只是一个静态近似没有删掉某个方向。

这个交互同时要求两种信息。$a_0$ 位于 $a_0$ 和 $a_1$ 的共同因果前缀，无论 $r$ 选择哪一个，$a_0$ 都先于 $r$；这种共同事实可以永久使用。但 $u$ 只先于晚发布，分析不能把两个来源身份合成一个事件，也不能把当前查询推得的 $u\mathrel\hb r$ 永久写入基础表示。

\subsection{固定观测与完整支持}
输入 $I$ 包含事件集合 $E$、线程顺序、事件类型与固定地址、写值、初始化、观测 $O$，以及可检查的程序位置和协议元数据 $\Lambda$。令 $n=|E|$，线程数为 $p$，地址数为 $a$。每条线程记录是所声明程序的无缺口有限前缀，允许在调用内部或等待中停止，全部失败尝试保留。$\Lambda$ 指明线程角色、协议地址所有权、调用位置、容量或参与者数及必要的布局约定；它在选择 RF／CO 前给出。本文不从任意无注释日志自动发现协议。

我们区分两类输入精度。完整值观测记录每次读值及已发生分支。分支观测允许某些返回值只由允许集合或已发生守卫约束。例如缓冲成功消费只说明两个索引不同，不能把未记录的对端返回值补成某个版本。写的具体地址与写值在两种模式中均已由固定事件前缀确定。

在任何协议过滤之前建立原始变量集合 $X$。每个读与所有满足其原始值观测的同地址写之间注册一个 RF 位；未记录值的读包含所有同地址写。对每个同地址无序写对注册一个 CO 方向位，反方向用该位的否定表达。初始化相关位也在内。后续证明某个位恒真或恒假时，只记录其常量，不删除或重编号。令 $\Phi_I(X)$ 为“恰一来源、逐地址初始化最小的全修改序、全部值／分支观测及式~\eqref{eq:ra}”的布尔函数。

\begin{definition}[条件完整支持]\label{def:support}
对原始文字合取 $\alpha\subseteq\Lit(X)$，若 $\Phi_I\land\alpha$ 可满足，则
\[
 \Supp_I(\alpha)=\{\lambda\in\Lit(X):\exists x.\ \Phi_I(x)\land\alpha(x)\land\lambda(x)\}.
\]
若不可满足，接口单独返回 $\UNSAT$。在可满足时，$\lambda$ 必然成立等价于其反文字不在支持集合中。
\end{definition}

返回集合涵盖全部原始 RF 和 CO 位，含已退休事件及永久常量。目标是对同一个固定 $I$ 反复增加和撤销 $\alpha$，不因条件变化而丢失原来可能的完成。

\subsection{先固定结构，再接受查询}
构造分两部分进行。协议分析从输入中推出可靠来源域 $D$、共同因果下界 $H$ 和共同写序下界 $L$，生成证书 $\Gamma_I$；在这些候选上建立 $G_I$ 及处理计划 $\pi_I$。随后精确编译得到 DAG $\DAG_I$，由它回答任意 $\alpha$。顺序可概括为
\begin{equation}\label{eq:quantifier}
 \exists(\Gamma_I,G_I,\pi_I)\quad
 \forall J\text{ 合法完成}\quad
 \exists\gamma_J:\ J\text{ 可由同一计划表示}.
\end{equation}
证明注释 $\gamma_J$ 可以记录该完成选了哪次发布及未记录的实际读值；调用编号、失败串位置和计划不能随 $J$ 改变。后续投影定理还会排除表示中的非法完成，因而比式~\eqref{eq:quantifier} 的覆盖条件更强。

共同前序是这一步的关键。例如 $a\mathrel\SB b$，且一个读可选同值来源 $a$ 或 $b$，虽然没有必选 RF，两个完成中都有 $a\mathrel\hb r$。只把必选 RF 加入 $\SB$ 再闭包会漏掉它。第~\ref{sec:frontier}节从这样的可靠关系定义未决因果前沿，并证明其余仍需连接未来的稳定写每地址至多保留 $p$ 个。第~\ref{sec:protocols}节解释如何由协议获得足够强的关系与连续计划；第~\ref{sec:compilation}节则证明关系状态的压缩不影响原始条件支持。
